\section{AlexNet：深度学习的引爆点}

\subsection{AlexNet 的历史意义}

\begin{figure}[H]
\centering
\includegraphics[width=0.95\textwidth]{slides-images/slide-063.jpg}
\caption{AlexNet 架构（2012）：深度学习革命的起点\protect\footnotemark}
\end{figure}
\footnotetext{Slides 第64页。}

AlexNet（Krizhevsky, Sutskever, Hinton, 2012）是在 ImageNet 大规模视觉识别挑战赛（ILSVRC 2012）上取得突破性成绩的卷积神经网络。它在 Top-5 错误率上以 $\sim$10\% 的优势击败了传统方法（从 $\sim$26\% 降到 $\sim$16\%），开启了深度学习在计算机视觉中的统治时代。

AlexNet 的成功因素：
\begin{itemize}
    \item 使用\textbf{GPU} 训练（当时是两块 GTX 580）
    \item 大规模数据集 ImageNet（120 万张图片）
    \item ReLU 激活函数（取代 sigmoid/tanh）
    \item Dropout 正则化
\end{itemize}

\begin{figure}[H]
\centering
\includegraphics[width=0.95\textwidth]{slides-images/slide-065.jpg}
\caption{ImageNet 竞赛中各年份最优模型的错误率和层数变化\protect\footnotemark}
\end{figure}
\footnotetext{Slides 第66页。}

\subsection{本章小结}

AlexNet 的架构本身并不复杂——它基本上就是 LeNet-5 的放大版本。它的革命性在于证明了三件事：（1）CNN 可以在大规模数据集上训练，（2）GPU 是深度学习的关键加速器，（3）端到端学习的特征远优于手工特征。

%% ========== 总结与延伸 ==========

